%%%%%%%%%%%%%%%%%%%%%%%%%%%%%%%%%%%%%%%%%%%%%%%%%%%%%%%%%%%%%%%%%%%%
%% I, the copyright holder of this work, release this work into the
%% public domain. This applies worldwide. In some countries this may
%% not be legally possible; if so: I grant anyone the right to use
%% this work for any purpose, without any conditions, unless such
%% conditions are required by law.
%%%%%%%%%%%%%%%%%%%%%%%%%%%%%%%%%%%%%%%%%%%%%%%%%%%%%%%%%%%%%%%%%%%%

\documentclass[
  digital, %% This option enables the default options for the
           %% digital version of a document. Replace with `printed`
           %% to enable the default options for the printed version
           %% of a document.
  twoside, %% This option enables double-sided typesetting. Use at
           %% least 120 g/m² paper to prevent show-through. Replace
           %% with `oneside` to use one-sided typesetting; use only
           %% if you don’t have access to a double-sided printer,
           %% or if one-sided typesetting is a formal requirement
           %% at your faculty.
  table,   %% This option causes the coloring of tables. Replace
           %% with `notable` to restore plain LaTeX tables.
  lof,     %% This option prints the List of Figures. Replace with
           %% `nolof` to hide the List of Figures.
  lot,     %% This option prints the List of Tables. Replace with
           %% `nolot` to hide the List of Tables.
  %% More options are listed in the user guide at
  %% <http://mirrors.ctan.org/macros/latex/contrib/fithesis/guide/mu/fi.pdf>.
]{fithesis3}
%% The following section sets up the locales used in the thesis.
\usepackage[resetfonts]{cmap} %% We need to load the T2A font encoding
\usepackage[T1,T2A]{fontenc}  %% to use the Cyrillic fonts with Russian texts.
\usepackage[
  main=english, %% By using `czech` or `slovak` as the main locale
                %% instead of `english`, you can typeset the thesis
                %% in either Czech or Slovak, respectively.
  english, german, czech, slovak %% The additional keys allow
]{babel}        %% foreign texts to be typeset as follows:
%%
%%   \begin{otherlanguage}{german}  ... \end{otherlanguage}
%%   \begin{otherlanguage}{russian} ... \end{otherlanguage}
%%   \begin{otherlanguage}{czech}   ... \end{otherlanguage}
%%   \begin{otherlanguage}{slovak}  ... \end{otherlanguage}
%%
%% For non-Latin scripts, it may be necessary to load additional
%% fonts:
\usepackage{paratype}
\def\textrussian#1{{\usefont{T2A}{PTSerif-TLF}{m}{rm}#1}}
%%
%% The following section sets up the metadata of the thesis.


% TODO packages
% https://tex.stackexchange.com/questions/9796/how-to-add-todo-notes
\usepackage{xargs}                      % Use more than one optional parameter in a new commands
\usepackage[pdftex,dvipsnames]{xcolor}  % Colored text etc.
\usepackage[colorinlistoftodos,prependcaption,textsize=tiny]{todonotes}
\newcommandx{\unsure}[2][1=]{\todo[linecolor=red,backgroundcolor=red!25,bordercolor=red,#1]{#2}}
\newcommandx{\change}[2][1=]{\todo[linecolor=blue,backgroundcolor=blue!25,bordercolor=blue,#1]{#2}}
\newcommandx{\info}[2][1=]{\todo[linecolor=OliveGreen,backgroundcolor=OliveGreen!25,bordercolor=OliveGreen,#1]{#2}}
\newcommandx{\improvement}[2][1=]{\todo[linecolor=Plum,backgroundcolor=Plum!25,bordercolor=Plum,#1]{#2}}
\newcommandx{\thiswillnotshow}[2][1=]{\todo[disable,#1]{#2}}

\usepackage{graphicx}
\usepackage{amsmath, amssymb, mathtools}
\usepackage{caption}
\usepackage{subcaption}
\usepackage{url}
\usepackage{siunitx}
\usepackage{ifthen}

\newcounter{camera_ready}  % camera ready switch, using heavy graphs, in PDF
\setcounter{camera_ready}{0}  % set to 1 to print in vectors
\newcommand{\ifcamready}[2]{\ifthenelse{\value{camera_ready}>0}{#1}{#2}}
\newcommand{\itembf}[1]{\item {\bf{#1}}}
\newcommand{\bigO}[0]{\mathcal{O}}

\thesissetup{
    date          = \the\year/09/05,
    university    = mu,
    faculty       = fi,
    type          = prop,
    author        = {Mgr. Dušan Klinec},
    gender        = m,
    advisor       = {prof. RNDr. Václav Matyáš, M.Sc., Ph.D.},
    title         = {Secure multi-party computation with resource-constrained devices},
    TeXtitle      = {Secure multi-party computation with resource-constrained devices},
    keywords      = {multi-party computation, secure hardware, cryptocurrencies},
    TeXkeywords   = {multi-party computation, tamper-proof tokens, multi signatures, smart cards, resource constrained devices, cryptocurrencies, monero, \ldots},
    abstract      = {This is the abstract of my thesis, which can

                     span multiple paragraphs.},
    thanks        = {These are the acknowledgements for my thesis, which can

                     span multiple paragraphs.},
    bib           = example.bib,
}
\usepackage{makeidx}      %% The `makeidx` package contains
\makeindex                %% helper commands for index typesetting.
%% These additional packages are used within the document:
\usepackage{paralist} %% Compact list environments
\usepackage{amsmath}  %% Mathematics
\usepackage{amsthm}
\usepackage{amsfonts}
\usepackage{framed}
\usepackage{color}
\usepackage{url}      %% Hyperlinks
\usepackage{markdown} %% Lightweight markup
\usepackage{listings} %% Source code highlighting
\lstset{
  basicstyle      = \ttfamily,%
  identifierstyle = \color{black},%
  keywordstyle    = \color{blue},%
  keywordstyle    = {[2]\color{cyan}},%
  keywordstyle    = {[3]\color{olive}},%
  stringstyle     = \color{teal},%
  commentstyle    = \itshape\color{magenta}}
\usepackage{floatrow} %% Putting captions above tables
\floatsetup[table]{capposition=top}
\definecolor{shadecolor}{rgb}{1,0.8,0.3}
\begin{document}

%\chapter*{Introduction}
%\addcontentsline{toc}{chapter}{Introduction}

%%%%%%%%%%%%%%%%%%%%%%%%%%%%%%%%%%%%%%%%%%%%%%%%%%%%%%%%%%%%%%%%%%%%%%%%%%%%%%%%%%%%%%%%%%%

\section{Introduction}
\begin{shaded}
\begin{itemize}
    \item Introduction, motivation
    \item Current state
    \item Problems, solved parts
    \item Open problems, hints what to do
    \item Aim of the work, high level.
    \item Chapters, contents.
    \item 2-3 pages
\end{itemize}
\end{shaded}

Itemized:
\begin{itemize}
	\item Bug resilience, Kleptography, Single point of failure of resource constrained devices. GACR proposal motivation. eID system failure due to {\bf{ROCA}}~\cite{2017-ccs-nemec}. Smart-ID~\cite{smart_id_ee}. 
    
    - TODO: ROCA 2 paper.
        
    \item {\bf{Fingerprinting}} public keys from cryptographic libraries and devices possible. Enables to target attack on particular vulnerable library. {\bf{ROCA}} also fingerprintable, like shown in our previous research {\it{Measuring Popularity of Cryptographic Libraries in Internet-Wide Scans}}~\cite{2017-acsac-nemec}.
    
        \item In previous research we devised powerful randomness test {\bf{\emph{Booltest}}} \cite{booltest_secrypt2017} which provides comparable and better results compared to standard randomness testing batteries. Testing randomness of PRNGs or cryptographic primitives is essential for the assessment of the quality of cryptographic schemes. It is usually the first step of a security audit. In the extended version of the paper submitted to the Springer Book of Secrypt We also tested output of TRNGs of smart cards and hardware wallets Trezor and Ledger. Having device failing to pass randomness test indicates a further problem with potentially high impact. Implementing the SMPC may help to avoid similar vulnerabilities and potential backdoors. 
    
    \item Many scenarios (e.g., DRM, licensing, cloud assisted computation) require computation with sensitive data on attacker computer. Universal approach is FHE but not practical. Another solution is a secure multi-party computation (SMPC), utilizing secure crypto processors (SCPs) or smart cards to provide better security guarantees. {\bf{Whitebox}} attack resistant cryptography \cite{whitebox_klinec_santacrypt2013,Klinec2013thesis} attempts to transform algorithms so their evaluation on attackers computer is secure, e.g. AES algorithm with embedded decryption key without possibility of key extraction. However, whitebox schemes often fail (cite) which indicates SMPC / SCP schemes are more practical.
    
    \item In previous research, we proposed and analyzed an Intrusion detection system running on resource-constrained devices creating wireless sensor networks \cite{wsnprotectlayer}. Effective multi-party computation algorithms could improve security properties of {\bf{WSN}} protocols.
    
    \item SMPC with Smartcards as our previous research A {\it {\bf{Touch of Evil}}: High-Assurance Cryptographic Hardware from Untrusted Components} \cite{2017-ccs-mavroudis}. Multi-Encryption and multi-signatures N-of-N. At least one honest participant is enough to protect secret.
    
    \item Paper with V. Sedlacek analyzing particular {\bf{number factorization}} method based on elliptic curves. Submitted by the end of theses submission deadline.

\end{itemize}

\subsection{Motivation - text}

    {\bf{Mix of GACR and Monero paper. Broader.}}
    
    The focus of the research aimed by theses is to analyze a secure distributed function computation in restricted environments. The computation model consists of multiple heterogeneous nodes. Node is either a powerful generic device such as cloud computation servers or very resource-restricted device, e.g., with low computational power, memory, energy. The typical instantiation of the resource-restricted device is Secure cryptographic processors (SCP) or smart card with additional properties such as tamper protection. 
    We aim to study ECC based protocols enabling a secure computation on powerful but potentially malicious devices with the help of trusted SCPs and to improve the system robustness if the SCP is compromised.
    
    \paragraph{SCPs} play a crucial role in many critical systems across various fields, e.g., banking, military, and government.
    SCPs provide an important additional layer of protection against attacks. A potential security vulnerability or a bug in the SCP can have broad consequences and compromise overall system security. For example, our recent discovery of a security vulnerability in RSA key generation in the Infineon SCP shows the fragility of eID systems having the affected chip a single point of failure \cite{2017-ccs-nemec}. 
    
    Therefore it is highly important to test for potential
    bugs and backdoors in the SCPs as critical system components. A manufacturing path from blueprints to a final product has multiple steps, and the manufacturer is not able to oversee all parts of the process. Backdoors can be induced at any stage, and even product certification cannot give reasonable assurances for critical systems \cite{2017-ccs-nemec}. Moreover the Kleptography-based \cite{Young:1997:KUC:1754542.1754551} backdoor can be present in the standard itself so only the attacker can detect and use it while for others it is computationally indistinguishable, e.g., RSA key generation. 
    
    The solution we propose to this family of threats is removing the single point of failure by employing a secure multi-party computation (SMPC). 
    With the SMPC multiple independent participants in the protocol can jointly compute an agreed-upon function so that individual function inputs of the participants remain private while the output is computed by all participants. The participants can be malicious, deviate from the protocol, and collude to get the private function inputs of others. SMPC is also known as secure function evaluation (SFE) or privacy-preserving computation. Several SMPC protocol exists with different guarantees with respect to the attacker model, a number of colluding participants and so on.
    
    \paragraph{Smart-ID} A practical example of SMPC deployed in Estonia is the Smart-ID \cite{smart_id_ee}, used for authentication and legal digital signatures. Smart-ID extends the RSA digital signature scheme, so it is computed jointly by two independent parties, user's smart-phone, and the server-side HSM. 
    Each party creates a signature share $s_i$ using a private key share $p_i$. The final composite signature $s$ is a simple combination of the signature shares $s_i$. 
    The $s$ is RFC3447 \cite{rfc3447} RSASSA-PKCS1-v1.5 compliant signature under the composite private key $p$. Even after compromising one of the $p_i$ the attacker is not able to forge signatures.


%%%%%%%%%%%%%%%%%%%%%%%%%%%%%%%%%%%%%%%%%%%%%%%%%%%%%%%%%%%%%%%%%%%%%%%%%%%%%%%%%%%%%%%%%%%
%%%%%%%%%%%%%%%%%%%%%%%%%%%%%%%%%%%%%%%%%%%%%%%%%%%%%%%%%%%%%%%%%%%%%%%%%%%%%%%%%%%%%%%%%%%

\section{State of the Art}
\begin{shaded}
\begin{itemize}
    \item More concrete introduction
    \item Current state of the art
    \item Classical and up to date results, comparison, overview / survey
    \item Up to date topics to address, who is solving those
    \item How to contribute
    \item Range: 8-12 pages
\end{itemize}
\end{shaded}


\section{Related work - itemized}
\begin{itemize}
    \item {\bf{SMPC}}, multi party computation in general. Secure function evaluation, privacy preserving computation, specific algorithms, e.g., privacy preserving set intersection, database lookup, data mining. 
    
    \item {\bf{Fully homomorphic encryption}}, alternative solution, not practical for universal problems. Could substitute need for secure processors. Move to bottom, like SoK paper, just mention.
    
    \item {\bf{Secure processor assisted cryptography}}, using FPGAs, smart-cards, 
    
    \item {\bf{Multi-signatures}} multiparty signature schemes. 1-N, N-N, k-N. Popular in cryptocurrencies.
    
    \item {\bf{Privacy centric cryptocurrencies}}. The initial goal is to support transaction signing on resource-constrained devices. Key for spending funds is a highly valuable target for attackers, sometimes worth millions of dollars. Storing key in a secure processor lowers the attack surface compared to general PC wallet. SCP can provide further security guarantees, security proofs / high EAL certification. 
    
    Privacy-centric cryptocurrencies use complex crypto schemes which makes support for SCP difficult due to the restricted computational environment. HW wallet for Monero, ZCoin, Zcash or derivation of those. 
    
    Another goal is multi-signature support, so one signing key is distributed among more hardware wallets from different vendors. Secure blockchain scanning is another topic. 
\end{itemize}

\section{Related work - DRAFT}
    \subsection{SMPC}
    - garbled circuits, ORAMs, circuit compilers.
    
    \subsection{Fully homomorphic encryption}
    - gentry, ideal solution to trusted computing on cloud or with secret inputs. not practical for general problems.
    
    \subsection{Multi signatures}
    - schnorr, mention UCL paper, okamoto, public key aggregation,  
    
    \subsection{Secure processor assisted cryptography}
    - roca, add DH proposal, FPGAs, 
    
    \subsection{Privacy centric cryptocurrencies}
    - multisigs, monero, zcoin, zcash
    
\newpage    
\section{Related work}

Secure multi-party computation (MPC) is very generic and well studied cryptographic task.
It allows mutual distrustful parties $P_i, i \in \{1,\dots,m\}$ to jointly perform a distributed computation of a function $f(x_i,\dots,x_m) = (y_i,\dots,y_m)$ on their inputs $x_i$ in a way each party learns only $y_i$ from the execution (\emph{privacy}). In particular, other users' inputs $x_{j, j \neq i}$ and outputs $y_{j, j \neq i}$ remain secret for user $P_i$. Depending on the security properties of the particular MPC protocols they should remain secure even with presence of a coalition of colluding corrupted nodes. MPC basically attempts to emulate a trusted third party (TTP) computing the function $f$ on user's behalf preserving the \emph{correctness} and \emph{privacy} as described.

MPC protocols are characterized with respect to various criteria:
\begin{itemize}
    \itembf{Number of parties.} Two-party computation is an important special case $m=2$ which has dedicated body of research as it achieves own specific feasibility results due to nature of the problem, i.e., there is no honest majority and often motivated by popular client-server paradigm.

    \itembf{Attacker types.} The \emph{semi-honest} attackers may form a colluding coalitions with the aim to learn more information from the protocol runs than the privacy property of the MPC describes. They may communicate and share information from runs. However, they exactly follow the protocol definition. This attacker model is also called \emph{passive} or \emph{honest-but-curious}. 
    On the other hand, \emph{malicious} or \emph{active} attackers may on top of that perform arbitrary actions during the protocol execution. They can diverge from the protocol, forge, replay messages, stop responding etc.

    \emph{Static} attackers corrupt the nodes on the beginning of the protocol. Corrupted set is unknown to the honest parties. \emph{Adaptive} attackers may selectively choose which node to corrupt during the protocol execution based on the collected information from the run. Intuitively, adaptive attacker model is stronger.
    
    \itembf{Size of the attacker coalition.} Basic security property of MPC protocol is a size of the coalition $t$ the MPC protocol stays secure. Depending on the attacker model it can be \emph{dishonest majority} up to $m-1$, \emph{honest majority} $2t+1 < m$, which makes sense only with $m>2$. There are also special cases requiring $2/3$ majority.

    \itembf{Computational power.} The adversary is typically assumed to be \emph{computationally bounded}, i.e., it's running time is probabilistic polynomial in the inputs $x_i$ and number of parties $m$.
    MC protocols with this attacker typically rely on typical cryptographic tools and hardness assumptions like Decisional Diffie-Hellman assumption, factoring hardness assumption and so on. From this reason the model is also called \emph{cryptographic}.
    
    Contrary to that \emph{information-theoretic} model poses no restrictions on attackers computational power. It's also known as \emph{computationally unbounded model} or \emph{perfectly private}. MPC in this model provides higher security guarantees but the security breaks completely if the attacker coalition size $t$ grows above the required limit. MPC in this model are typically using perfectly-secure one-time-pad and secret sharing, e.g., Shamir's secret sharing. They are usually computationally more effective as there is no cryptography related but the communication cost may render them ineffective. 
    
    \itembf{Set-up assumptions.} MPC security is stated with respect to various set-up assumptions. If not stated the MPC is assumed to work in so called \emph{plain model}. It may also use stronger security models such as \emph{Common Reference String (CRS)} model assuming all nodes have access to a CRS from certain distribution generated by a TTP. It may be e.g., cryptographic keys or randomness used later in the protocol. Other commonly used models are \emph{Correlated Randomness Model} and \emph{Public key infrastructure (PKI)}. 
    
    Orthogonal to that are theoretical assumptions such as \emph{One way functions}, \emph{Random oracle model}, \emph{collision resistant hash functions}.
    
    \itembf{Communication properties.} MPC typically assume \emph{private channels}, i.e., private and reliable point-to-point communication channels. Some properties may additionally require a \emph{broadcast}. Channels can be either \emph{synchronized} or \emph{asynchronous}. 
    
    Protocols are usually characterized with the respect to communication \emph{rounds}. In one round each participant $P_i$ may send message to others $\overline{P_i} = \{P_j; i \neq i\}$. Some literature define round as one message round-trip, i.e., message exchange. Literature agrees that if a message is a reaction on previous message it is considered as a new round. 
    
    With respect to number of rounds are protocols characterized as \emph{polynomialy input or function bound}, \emph{constant-round}.
\end{itemize}

\paragraph{Two-party computation.} %
One of the founding work on MPC is a seminal paper of Yao \cite{Yao86} which proved \emph{completeness} of the two-party computation, i.e. all computable function $f$ can be computed in secure two-party computation model in semi-honest setting with computationally bounded attackers. 

Yao's approach is based on a converting function $f$ to a boolean circuit $C$. The $C$ consists of gates, i.e., nodes performing addition and multiplication (AND and XOR) on its inputs and producing an output. Number of input wires on gate is fan-in, typically 2 and number of outputs is called fan-out which is typically 1. Gates are connected with wires into a circuit. Input wires are connecting user inputs to gate's input, output wires are connected to gates outputs and unconnected on the other end.

Circuit \emph{creator} creates the circuit $C$ and \emph{garbles} it in a sense he encrypts and permutes the gates truth tables to protect his values and intermediate state. This operation gave the scheme name \emph{garbled circuits}. Creator sends the garbled circuit with his encrypted inputs, i.e., wire \emph{labels}, to the \emph{evaluator}. In order to evaluate the circuit she needs to know wire labels corresponding to her inputs and only the creator knows the encryption key. The $\binom{2}{1}$ Oblivious Transfer is used to query corresponding values from the creator without revealing evaluator's inputs. Evaluator evaluates the $C$ in topological order gate-by-gate and obtains output labels which are decrypted by the creator to obtain final result of the computation. The protocol is constant round, independent on the circuit depth. 
Obviously the scheme works only with semi-honest attackers as the malicious creator could create different circuit $C^\prime$ leaking the evaluator's inputs. The more formal description can be found in \cite{BHR12}. 

\paragraph{Multi-party extension.} %
Yao's work has been extended by \cite{GMW87} to a multi-party case and with malicious adaptive nodes using only one-way functions. They make use of Verifiable Secret Sharing (VSS) and agreement protocols. The work is based on boolean circuit but compared to \cite{Yao86} each party evaluates the circuit independently with shares of other players inputs. This work is a completeness theorem for distributed protocols with honest majority. Two protocols were proposed, the first for semi-honest setting which is secure with $m-1$ corrupted nodes, the other in the malicious settings secure with honest majority (correctness and privacy). Later \cite{CLOS02} shown how to construct a multi-party protocols in universally composable framework in computational assumptions with dishonest majority up to $m-1$ of adaptive malicious parties. 

\paragraph{Computationally unbounded.}
The information-theoretic equivalent of this work was done by \cite{BGW88} and \cite{CCD88}. The work is based on Shamir's secret sharing and arithmetic circuits in finite fields. The malicious case make use of error correcting Reed Solomon codes. They proved a honest majority is required for semi-honest setting in this model and $2/3$-honest majority for malicious settings. Moreover, these bounds are tight. These protocols have $\bigO(\text{depth}(C))$ communication rounds.

\paragraph{Extensions.} A large body of research has been based on \cite{Yao86, GMW87, BGW88, CCD88} and many extensions and improvements were published. 
For the garbled circuits it entails: Point-and-permute \cite{BMR90}, Free XOR \cite{KS08}, garbled row reduction \cite{NPS99, PSSW09}, FleXOR \cite{KMR14}, half gates \cite{ZRE14}, garbled gadgets \cite{BMR16} and CompGC \cite{cryptoeprint:2016:458}.

Cut-and-choose extension \cite{LP07,LP11,Lin16} enables use of the garbled circuit with malicious parties. The approach has been extensively studied and significantly improved, e.g., LEGO \cite{NO09} and it's numerous improvements \cite{cryptoeprint:2013:155, FJNT15, cryptoeprint:2016:1069, cryptoeprint:2017:226, KNBTR17, cryptoeprint:2018:578}.

Secret-sharing based approach extensions working with malicious parties: SPDZ \cite{DPSZ12, cryptoeprint:2012:642}, MiniMac \cite{DZ13}, TinyOT \cite{NNOB12}, MASCOT \cite{KOS16}, TinyTable \cite{cryptoeprint:2016:695}, and further improvements \cite{FKOS15, DZ16, LPJ17, BCGIM17, cryptoeprint:2017:1230}. The \cite{WRK17} generalizes TinyOT to multiparty setting, dishonest majority of malicious attackers in standalone security model while \cite{DPSZ12} is statistically UC-secure.

\paragraph{Constant-round protocols.} The \cite{BB89} was the first Constant-round information-theoretic protocol. Later \cite{IK00, IK02} used so called randomized polynomials and randomized encodings \cite{AIK04} of a function to construct constant round protocols in information-theoretic setting with active adversary. Later \cite{ABT18} proved a tight bound 2-rounds for perfectly secure semi-honest MPC with honest majority using a new tool multi-party randomized encodings. The \cite{cryptoeprint:2018:1078} shown 2-round information-theoretic statistically secure MPC with malicious adversaries and honest majority with selective abort which was shown by \cite{GIKR04} to be the strongest notion of security for 2-round protocols.

\paragraph{Universally composable framework.}%
Nananana nana nan anana


\newpage
\begin{itemize}
    \item UC secure, setup assumptions, tokenzzzz
\end{itemize}

UC-related: "Consequently, many known protocols (e.g., the general protocol of
[gmw87], to name one) are not UC secure. Thus, the feasibility of realizing cryptographic tasks
requires re-investigation within the UC framework. Let us briefly summarize the known results.
In the case of a majority of honest parties, there exist UC secure protocols for computing any
functionality ([c01], building on [bgw88, rb89, cfgn96]). Also, in the honest-but-curious case
(i.e., when even corrupted parties follow the protocol specification), UC secure protocols exist for
essentially any functionality [clos02]."



Secure function evaluation has been broadly studied \cite{
yao82, conf/sp/SonghoriHS0K15, Malkhi:2004:FST:1251375.1251395, 10.1007/978-3-642-40041-4_5}, namely with respect to private database search \cite{Bertino2005, Naor:1999:PPA:336992.337028}, private data mining \cite{Aldeen2015} and voting. However, it typically suffers from large performance or communication overhead.
The previous research uses FPGA as a trusted element, trading trust for computational efficiency. However, FPGA itself cannot withstand attacks by side channel power analysis or fault induction attacks.
SCPs, on the other hand, are tamper-proof but do not offer the flexibility of FPGAs and, namely, they suffer from low computational power and very limited memory \cite{10.1007/978-3-642-19074-2_1, Boureanu:2015:LCC:2714576.2714591, cryptoeprint:2016:615}.
In \cite{2017-ccs-mavroudis}, we designed, theoretically analyzed, and evaluated \emph{Myst}, a suite of ECC-based cryptographic protocols for a distributed key generation, encryption, decryption and signature generation using SMPC on smartcards.
Also, the K-out-of-N ECC-based multi-signature protocol research has recently undergone rapid development, especially due to rising crypto-currencies popularity as they rely on multi-signatures for signing outgoing transactions in accounts held by multiple owners ~\cite{Bellare:2007:UAS:2394539.2394589, Bnz2017BulletproofsSP, cryptoeprint:2018:068}. 


%%%%%%%%%%%%%%%%%%%%%%%%%%%%%%%%%%%%%%%%%%%%%%%%%%%%%%%%%%%%%%%%%%%%%%%%%%%%%%%%%%%%%%%%%%%


\section{Thesis aim}
\begin{shaded}
\begin{itemize}
    \item What new can be expected
    \item What will be created
    \item Hypotheses to research
    \item Methods used
    \item Plans, time harmonogram
    \item Conferences
    \item Range: 2 pages
\end{itemize}
\end{shaded}

chapter{Study plan}

\begin{itemize}
	\item Transaction signing for privacy centric cryptocurrencies on secure crypto processors.  
    \item Multisignatures for privacy centric cryptocurrencies on secure crypto processors. 
    \item Multisig generalization k-of-N for privacy centric cryptocurrencies on secure crypto processors. 
    \item Analysis of existing protocols using SMPC. 
\end{itemize}

I will study and design multi-party protocols with different levels of trust and computational power of interacting parties, enabling a secure computation on powerful yet potentially malicious devices with the help of trusted SCPs.

SMPC-based techniques for systematic protection from the vulnerabilities and backdoors in ECC-based systems will be devised. These are aimed to remain secure even after the compromise of several cooperating parties. The primary focus is on SMPC ECC signature schemes in resource-constrained environments.

I believe that an exploration of a currently under-researched combination of highly secure yet very resource-limited devices and multi-party computation will bring significant benefits. The main benefit is the removal of a single point of failure from critical large-scale security systems, such as eID. The research will benefit from our expertise in this field and our previous results (incl. the Infineon cryptographic vulnerability analysis) \cite{2017-ccs-mavroudis, 2017-ccs-nemec}, multi-party protocol framework \emph{Myst}.

\paragraph{K-of-N MPC generalization.} In \cite{2017-ccs-mavroudis} we designed, theoretically analyzed, and evaluated the \emph{Myst}, suite of cryptographic protocols for a distributed key generation, encryption, decryption and signature generation using multi-party computation on smart cards. The protocols type is N-of-N, so an attacker has to compromise all involved components to gain any advantage. N-of-N increases resilience against the single point of failures and potential backdoors in the cards hardware (if selected from diverse manufacturers and supply chains). We would like to investigate options for generalization of these protocols to k-of-N threshold scheme yet still execute in a resource-constrained environment of the smart cards. This is an unsolved problem with broad potential applications, e.g., in the transaction signing in public blockchains as it helps to increase the availability and flexibility.  


\paragraph{Multi-signatures.} K-of-N ECC-based multi-signature protocol research has recently undergone a rapid development especially due to rising crypto-currencies popularity as they rely on multi-signatures for signing outgoing transactions in accounts held by multiple owners ~\cite{Bellare:2006:MPP:1180405.1180453, Bellare:2007:UAS:2394539.2394589, 10.1007/3-540-36288-6_3, cryptoeprint:2015:996, Bnz2017BulletproofsSP, cryptoeprint:2018:068}.
Recently a new multi-signature scheme \cite{cryptoeprint:2018:068} based on Schnorr signatures \cite{schnor-10.1007/3-540-46885-4_68} has been published with a strong potential in blockchain-based crypto-currencies due to its unique property set: 
a)~2-step protocol\footnote{Protocols with similar properties usually have three steps.}, b)~aggregated public key, c)~space effectiveness\footnote{All transaction inputs can be signed with one signature, compared to the current Bitcoin version where each input is signed separately.}. 
The key aggregation means the resulting signature is verifiable as the original Schnorr signature using a single aggregated public key which saves space. As the signature is saved in the blockchain the space savings are of high importance - the bigger the transaction, the more fees user need to pay for the transaction processing. This new perspective stresses the space effectiveness in the signature protocols design. We plan to investigate the possible space requirements reductions in the signature design as was performed in the case of Borromean signatures \cite{borromean} which reduced signature size almost in half. Moreover, we plan to study mechanisms of adapting signature schemes to the SCPs model. The motivation is a rising popularity of secure hardware, e.g., used to protect crypto-currency assets while signature protocols are not optimized for SCPs which brings challenges using SCPs in this ecosystem. The transformation is not straightforward and without protocol modifications.


The protocols will be analyzed and transformed into the form usable in the restricted environments such as SCPs. Results will be evaluated on the high-throughput smart card grid already used in \cite{2017-ccs-mavroudis}. 


Phase 2 focuses on a delegated transaction scanning and a secure protocol offloading from SCPs to hosts.
 
There are several general purposes, universal SMPC protocols which often suffer high computation and message complexity as a price for the universality. On the other hand, there are few very specific optimized SMPC solving particular problems which can be an order of magnitudes faster than universal protocols. Two-party secure computation is another specification of the universal model allowing better performance with a fixed number of participants \cite{Wang:2017:AGE:3133956.3134053}. We aim to study specialized ECC-based SMPC to address open problems in cryptocurrency model and to enable computation offloading to powerful hosts with use of SCPs.
%\todo{More citations on general and specific MPCs}
%\todo{Rethink if still actual: We need to add more generic basic research motivation with text below as practical demonstration of one potential use case}

\paragraph{SMPC in cryptocurrencies.} 
Specialized SMPC protocols can address novel challenges in the cryptocurrency model systems. There are several cryptocurrency-related open problems where users need to compute a joint function without disclosing their inputs.

Privacy-centric cryptographic currencies (e.g., Monero, Zcash, Dash, PivX\cite{monero_1098}) provide transaction unlinkability and privacy, usually implemented by defining two individual keys instead of just one: view-key and spend-key. A user scans the whole blockchain and performs a simple operation with his view-key on a per-transaction basis to determine whether the transaction being inspected is destined to him and to reveal the amount being sent. The spend-key is then used to spend the received money.
The view-key can be passed to another third party auditor as a read-only key to verify incoming transactions when needed (e.g., regulatory demands). An obvious downside of passing the raw view-key is the third party can leak it or use it to scan the blockchain even after the required auditory period. 
 
We plan to investigate multi-party protocols for a secure and controlled view-key sharing with another party. We plan to adapt existing SMPC protocols to keep the view-key private and fulfill regulatory demands of the authorities by actively participating in the verification process on a per-transaction basis. The benefit is that after the requirement invalidation the authority cannot use the view-key share to gain any advantage to further scan for a transaction without the owner's knowledge. The participation of the owning entity is required.

The increasing demand and usage of SCP as hardware wallets motivates the design and analysis of protocols using and/or supporting SCPs. Ideally, the view-key and spend-keys are stored in the SCP and never leave it. As the SCP is a low-performance device, it is challenging to scan the blockchain on the device directly. The same holds for participation in the SMPs.

\paragraph{Secure computation offloading.}
Secure computation offloading to the cloud with potentially malicious computation nodes is still an attractive open problem. There are several general solutions such as Fully Homomorphic encryption and Garbled circuits which are either universal and very ineffective or very specific. These mechanisms also often suffer high communication complexity.\todo{FHE, Garbled citations} 
We focus on asymmetric SMPC with resource-restricted devices (SCPs) in a sense the majority of the computation is performed on the powerful yet malicious host without compromising the security of the private function inputs. SCPs typically provide sensitive inputs to the function being computed, e.g., private keys. 

In the crypto-currency setup the zero-knowledge proofs, commitments, multi-signatures, and blockchain scanning would significantly benefit from effective computation offloading and multi-party protocol versions.

We aim to investigate new possibilities of the secure view-key sharing in the SCP via SMPC by offloading costly operations to the high-performance hosts, minimizing the number of rounds, operations and memory in SCP.


%%%%%%%%%%%%%%%%%%%%%%%%%%%%%%%%%%%%%%%%%%%%%%%%%%%%%%%%%%%%%%%%%%%%%%%%%%%%%%%%%%%%%%%%%%%

\section{Achieved results}
\begin{shaded}
\begin{itemize}
    \item What we have, created
    \item Conferences presented
    \item Range: 3-5 pages
\end{itemize}
\end{shaded}

\subsection{Monero paper}

Idea: submit to Oakland (rolling deadline) or some other A, A+ conference before theses submission. 

Covering basic Monero transaction, transaction signing protocol for hardware wallets, with security proof (ideally. reduction. original signing algorithm is secure then transformed protocol is secure as well).

Covering also blockchain scanning, so view key is not leaked to the host. Schemes designed with their properties. Group payment identification pseudonymity, fast scanning. Encrypted Bloom filters. Impossibility on some attractive scheme properties (publicly derivable public key chain for forward secrecy, corresponding private key derivation only with hardware wallet. Formalize). 


\section{Author's Publications}
\begin{shaded}
\begin{itemize}
    \item Contribution to the paper, percentual contribution
\end{itemize}
\end{shaded}

%%%%%%%%%%%%%%%%%%%%%%%%%%%%%%%%%%%%%%%%%%%%%%%%%%%%%%%%%%%%%%%%%%%%%%%%%%%%%%%%%%%%%%%%%%%

%% Bibliography
  \printbibliography[heading=bibintoc] 


%% Index
%   \makeatletter\thesis@blocks@clear\makeatother
%   \phantomsection %% Print the index and insert it into the
%   \addcontentsline{toc}{chapter}{\indexname} %% table of contents.
%   \printindex


%% Appendix
\appendix 
\chapter{Attached papers}

\section{Monero paper}
Submitted (ideally). Main author, 70\%.

\section{A Touch of Evil: High-Assurance Cryptographic Hardware from Untrusted Components}
Primarily targeting topic, 5\% \cite{2017-ccs-mavroudis}

\section{The Return of Coppersmith's Attack: Practical Factorization of Widely Used RSA Moduli}
ROCA. Motivation for SMPC, 5-10\% \cite{2017-ccs-nemec}.
Add ROCA extension.

\section{The Efficient Randomness Testing using Boolean Functions}
Booltest. \cite{booltest_secrypt2017}.
Add Booltest extended paper in journal.

\section{Measuring Popularity of Cryptographic Libraries in Internet-Wide Scans}
Fingerprinting. Significant contribution 20-30 \% \cite{2017-acsac-nemec}

\section{WSNProtectLayer – security middleware for wireless sensor networks}
Minor contribution. Resource-restricted devices. \cite{wsnprotectlayer}

\section{Whitebox attack resistant cryptography}
\cite{whitebox_klinec_santacrypt2013, Klinec2013thesis} 

- Citations count 9: \url{https://scholar.google.com/scholar?cites=3284880228025768764&as_sdt=2005&sciodt=0,5&hl=sk}

\section{ECC Factorization paper}
Submitted, V. Sedlacek, small contribution, running experiments, results aggregation. 

\section{Papers in progress}
\begin{itemize}
	\item ROCA - extended attack. Returning information to the system.
    
	\item Automated Smart card APDU AFL assisted fuzzer with systematic power trace analysis feedback.
    
    \item Zero knowledge proof of secrets generation correctness for smart cards. Producing a proof that given public key was generated according to the prescribed standard and thus cannot contain a bug / kleptography backdoor. 
    
    \item Static Rabin-Miller basis on embedded cards, testing, detection, exploitation. Based on Bleichenbacher \cite{10.1007/978-3-540-30580-4_2}.
    
\end{itemize}
\end{document}



%%%%%%%%%%%%%%%%%%%%%%%%%%%%%%%%%%%%%%%%%%%%%
%% Trash
%%%%%%%%%%%%%%%%%%%%%%%%%%%%%%%%%%%%%%%%%%%%%



%Theses are rumoured to be the capstones of education, so I decided
%to write one of my own. If all goes well, I will soon have a
%diploma under my belt. Wish me luck!


%\begin{otherlanguage}{czech}
%czz
%\end{otherlanguage}

%\chapter{These are}
%\section{the available}
%\subsection{sectioning}
%\subsubsection{commands.}
%\paragraph{Paragraphs and}
%\subparagraph{subparagraphs are available as well.}
%Inside the text, you can also use unnumbered lists,
%\begin{itemize}
%  \item such as
%   \item this one
%   \begin{itemize}
%     \item     and they can be nested as well.
%     \item[>>] You can even turn the bullets into something fancier,
%     \item[\S] if you so desire.
%   \end{itemize}
% \end{itemize}
% Numbered lists are
% \begin{enumerate}
%   \item very
%   \begin{enumerate}
%     \item similar
%   \end{enumerate}
% \end{enumerate}
% and so are description lists:
% \begin{description}
%   \item[Description list]
%     A list of terms with a description of each term
% \end{description}
% The spacing of these lists is geared towards paragraphs of text.
% For lists of words and phrases, the \textsf{paralist} package
% offers commands
% \begin{compactitem}
%   \item that
%   \begin{compactitem}
%     \item are
%     \begin{compactitem}
%       \item better
%       \begin{compactitem}
%         \item suited
%       \end{compactitem}
%     \end{compactitem}
%   \end{compactitem}
% \end{compactitem}
% \begin{compactenum}
%   \item to
%   \begin{compactenum}
%     \item this
%     \begin{compactenum}
%       \item kind of
%       \begin{compactenum}
%         \item content.
%       \end{compactenum}
%     \end{compactenum}
%   \end{compactenum}
% \end{compactenum}
% The \textsf{amsthm} package provides the commands necessary for the
% typesetting of mathematical definitions, theorems, lemmas and
% proofs.

%% We will define several mathematical sectioning commands.
\newtheorem{theorem}{Theorem}[section] %% The numbering of theorems
                               %% will be reset after each section.
\newtheorem{lemma}[theorem]{Lemma}         %% The numbering of lemmas
\newtheorem{corollary}[theorem]{Corollary} %% and corollaries will
                               %% share the counter with theorems.
\theoremstyle{definition}
\newtheorem{definition}{Definition}
\theoremstyle{remark}
\newtheorem*{remark}{Remark}

% \begin{theorem}
%   This is a theorem that offers a profound insight into the
%   mathematical sectioning commands.
% \end{theorem}
% \begin{theorem}[Another theorem]
%   This is another theorem. Unlike the first one, this theorem has
%   been endowed with a name.
% \end{theorem}
% \begin{lemma}
%   Let us suppose that $x^2+y^2=z^2$. Then
%   \begin{equation}
%     \biggl\langle u\biggm|\sum_{i=1}^nF(e_i,v)e_i\biggr\rangle
%     =F\biggl(\sum_{i=1}^n\langle e_i|u\rangle e_i,v\biggr).
%   \end{equation}
% \end{lemma}
% \begin{proof}
%   $\nabla^2 f(x,y)=\frac{\partial^2f}{\partial x^2}+
%    \frac{\partial^2f}{\partial y^2}$.
% \end{proof}
% \begin{corollary}
%   This is a corollary.
% \end{corollary}
% \begin{remark}
%   This is a remark.
% \end{remark}

% \chapter{Floats and references}
% \begin{figure}
%   \begin{center}
%     %% PNG and JPG images can be inserted into the document as well,
%     %% but their resolution needs to be adequate. The minimum is
%     %% about 100 pixels per 1 centimeter or 300 pixels per 1 inch.
%     %% That means that a JPG or PNG image typeset at 4 × 4 cm should
%     %% be 400 × 400 px large at the bare minimum.
%     %%
%     %% The optimum is about 250 pixels per 1 centimeter or 600
%     %% pixels per 1 inch. That means that a JPG or PNG image typeset
%     %% at 4 × 4 cm should be 1000 × 1000 px large or larger.
%     \includegraphics[width=4cm]{fithesis/logo/mu/fithesis-base.pdf}
%   \end{center}
%   \caption{The logo of the Masaryk University at 40\,mm}
%   \label{fig:mulogo1}
% \end{figure}

% \begin{figure}
%   \begin{center}
%     \begin{minipage}{.66\textwidth}
%       \includegraphics[width=\textwidth]{fithesis/logo/mu/fithesis-base.pdf}
%     \end{minipage}
%     \begin{minipage}{.33\textwidth}
%       \includegraphics[width=\textwidth]{fithesis/logo/mu/fithesis-base.pdf} \\
%       \includegraphics[width=\textwidth]{fithesis/logo/mu/fithesis-base.pdf}
%     \end{minipage}
%   \end{center}
%   \caption{The logo of the Masaryk University at $\frac23$ and
%     $\frac13$ of text width}
%   \label{fig:mulogo2}
% \end{figure}

% \begin{table}
%   \begin{tabularx}{\textwidth}{lllX}
%     \toprule
%     Day & Min Temp & Max Temp & Summary \\
%     \midrule
%     Monday & $13^{\circ}\mathrm{C}$ & $21^\circ\mathrm{C}$ & A
%     clear day with low wind and no adverse current advisories. \\
%     Tuesday & $11^{\circ}\mathrm{C}$ & $17^\circ\mathrm{C}$ & A
%     trough of low pressure will come from the northwest. \\
%     Wednesday & $10^{\circ}\mathrm{C}$ &
%     $21^\circ\mathrm{C}$ & Rain will spread to all parts during the
%     morning. \\
%     \bottomrule
%   \end{tabularx}
%   \caption{A weather forecast}
%   \label{tab:weather}
% \end{table}

% The logo of the Masaryk University is shown in Figure
% \ref{fig:mulogo1} and Figure \ref{fig:mulogo2} at pages
% \pageref{fig:mulogo1} and \pageref{fig:mulogo2}. The weather
% forecast is shown in Table \ref{tab:weather} at page
% \pageref{tab:weather}. The following chapter is Chapter
% \ref{chap:matheq} and starts at page \pageref{chap:matheq}.
% Items \ref{item:star1}, \ref{item:star2}, and
% \ref{item:star3} are starred in the following list:
% \begin{compactenum}
%   \item some text
%   \item some other text
%   \item $\star$ \label{item:star1}
%   \begin{compactenum}
%     \item some text
%     \item $\star$ \label{item:star2}
%     \item some other text
%     \begin{compactenum}
%       \item some text
%       \item some other text
%       \item yet another piece of text
%       \item $\star$ \label{item:star3}
%     \end{compactenum}
%     \item yet another piece of text
%   \end{compactenum}
%   \item yet another piece of text
% \end{compactenum}
% If your reference points to a place that has not yet been typeset,
% the \verb"\ref" command will expand to \textbf{??} during the first
% run of
% \texttt{pdflatex \jobname.tex}
% and a second run is going to be needed for the references to
% resolve. With online services -- such as Overleaf -- this is
% performed automatically.

% \chapter{Mathematical equations}
% \label{chap:matheq}
% \TeX{} comes pre-packed with the ability to typeset inline
% equations, such as $\mathrm{e}^{ix}=\cos x+i\sin x$, and display
% equations, such as \[
%   \mathbf{A}^{-1} = \begin{bmatrix}
%   a & b \\ c & d \\
%   \end{bmatrix}^{-1} =
%   \frac{1}{\det(\mathbf{A})} \begin{bmatrix}
%   \,\,\,d & \!\!-b \\ -c & \,a \\
%   \end{bmatrix} =
%   \frac{1}{ad - bc} \begin{bmatrix}
%   \,\,\,d & \!\!-b \\ -c & \,a \\
%   \end{bmatrix}.
% \] \LaTeX{} defines the automatically numbered \texttt{equation}
% environment:
% \begin{equation}
%   \gamma Px = PAx = PAP^{-1}Px.
% \end{equation}
% The package \textsf{amsmath} provides several additional
% environments that can be used to typeset complex equations:
% \begin{enumerate}
%   \item An equation can be spread over multiple lines using the
%     \texttt{multline} environment:
%     \begin{multline}
%       a + b + c + d + e + f + b + c + d + e + f + b + c + d + e +
% f \\
%       + f + g + h + i + j + k + l + m + n + o + p + q
%     \end{multline}

%   \item Several aligned equations can be typeset using the
%     \texttt{align} environment:
%     \begin{align}
%               a + b &= c + d     \\
%                   u &= v + w + x \\[1ex]
%       i + j + k + l &= m
%     \end{align}

%   \item The \texttt{alignat} environment is similar to
%     \texttt{align}, but it doesn't insert horizontal spaces between
%     the individual columns:
%     \begin{alignat}{2}
%       a + b + c &+ d       &   &= 0 \\
%               e &+ f + g   &   &= 5
%     \end{alignat}

%   \item Much like chapter, sections, tables, figures, or list
%     items, equations -- such as \eqref{eq:first} and
%     \eqref{eq:mine} -- can also be labeled and referenced:
%     \begin{alignat}{4}
%       b_{11}x_1 &+ b_{12}x_2  &  &+ b_{13}x_3  &  &             &
%         &= y_1,                   \label{eq:first} \\
%       b_{21}x_1 &+ b_{22}x_2  &  &             &  &+ b_{24}x_4  &
%         &= y_2. \tag{My equation} \label{eq:mine}
%     \end{alignat}

%   \item The \texttt{gather} environment makes it possible to
%     typeset several equations without any alignment:
%     \begin{gather}
%       \psi = \psi\psi, \\
%       \eta = \eta\eta\eta\eta\eta\eta, \\
%       \theta = \theta.
%     \end{gather}

%   \item Several cases can be typeset using the \texttt{cases}
%     environment:
%     \begin{equation}
%       |y| = \begin{cases}
%         \phantom-y & \text{if }z\geq0, \\
%                 -y & \text{otherwise}.
%       \end{cases}
%     \end{equation}
% \end{enumerate}
% For the complete list of environments and commands, consult the
% \textsf{amsmath} package manual\footnote{
%   See \url{http://mirrors.ctan.org/macros/latex/required/amslatex/math/amsldoc.pdf}.
%   The \texttt{\textbackslash url} command is provided by the
%   package \textsf{url}.
% }.

% \chapter{\textnormal{We \textsf{have} \texttt{several} \textsc{fonts}
%   \textit{at} \textbf{disposal}}}
% The serified roman font is used for the main body of the text.
% \textit{Italics are typically used to denote emphasis or
% quotations.} \texttt{The teletype font is typically used for source
% code listings.} The \textbf{bold}, \textsc{small-caps} and
% \textsf{sans-serif} variants of the base roman font can be used to
% denote specific types of information.

% \tiny We \scriptsize can \footnotesize also \small change \normalsize
% the \large font \Large size, \LARGE although \huge it \Huge
% is \huge usually \LARGE not \Large necessary.\normalsize

% A wide variety of mathematical fonts is also available, such as: \[
%   \mathrm{ABC}, \mathcal{ABC}, \mathbf{ABC}, \mathsf{ABC},
%   \mathit{ABC}, \mathtt{ABC}
% \] By loading the \textsf{amsfonts} packages, several additional
% fonts will become available: \[
%   \mathfrak{ABC}, \mathbb{ABC}
% \] Many other mathematical fonts are available\footnote{
%   See \url{http://tex.stackexchange.com/a/58124/70941}.
% }.

% \chapter{Using lightweight markup}
% \shorthandoff{-}
% \begin{markdown*}{%
%   hybrid,
%   definitionLists,
%   footnotes,
%   inlineFootnotes,
%   hashEnumerators,
%   fencedCode,
%   citations,
%   citationNbsps,
% }

% If you decide that \LaTeX{} is too wordy for some parts of your
% document, there are [packages](https://www.ctan.org/pkg/markdown
% "Markdown") that allow you to use more lightweight markup next
% to it.

%  ![logo](fithesis/logo/mu/fithesis-base.pdf "The logo of the
%   Masaryk University")

% This is a bullet list. Unlike numbered lists, bulleted lists
% contain an **unordered** set of bullet points. When a bullet point
% contains multiple paragraphs, the list is typeset as follows:

%   * The first item of a bullet list

%     that spans several paragraphs,
%   * the second item of a bullet list,
%   * the third item of a bullet list.

% When none of the bullet points contains multiple paragraphs, the
% list has a more compact form:

%   * The first item of a bullet list,
%   * the second item of a bullet list,
%   * the third item of a bullet list.

% Unlike a bulleted list, a numbered list implies chronology or
% ordering of the bullet points. When a bullet point
% contains multiple paragraphs, the list is typeset as follows:

%   1. The first item of an ordered list

%      that spans several paragraphs,
%   2. the second item of an ordered list,
%   3. the third item of an ordered list.
%   #. If you are feeling lazy,
%   #. you can use hash enumerators as well.

% When none of the bullet points contains multiple paragraphs, the
% list has a more compact form:

%   6. The first item of an ordered list,
%   7. the second item of an ordered list,
%   8. the third item of an ordered list.

% Definition lists are used to provide definitions of terms. When
% a definition contains multiple paragraphs, the list is typeset
% as follows:

% Term 1

% :   Definition 1

% *Term 2*

% :   Definition 2

%         Some code, part of Definition 2

%     Third paragraph of Definition 2.

% When none of the bullet points contains multiple paragraphs, the
% list has a more compact form:

% Term 1
% :   Definition 1
% *Term 2*
% :   Definition 2

% Block quotations are used to include an excerpt from an external
% document in way that visually clearly separates the excerpt from
% the rest of the work:

% > This is the first level of quoting.
% >
% > > This is nested blockquote.
% >
% > Back to the first level.

% Footnotes are used to include additional information to the
% document that are not necessary for the understanding of the main
% text. Here is a footnote reference^[Here is the footnote.] and
% another.[^longnote]

% [^longnote]: Here's one with multiple blocks.

%     Subsequent paragraphs are indented to show that they
% belong to the previous footnote.

%         Some code

%     The whole paragraph can be indented, or just the first
%     line.  In this way, multi-paragraph footnotes work like
%     multi-paragraph list items.

% Citations are used to provide bibliographical references to other
% documents. This is a regular citation~[@borgman03, p. 123]. This is
% an in-text citation: @borgman03\. You can also cite several authors
% at once using both regular~[see @borgman03, p. 123; @greenberg98,
% sec.  3.2; and @thanh01] and in-text citations: @borgman03 [p.123;
% @greenberg98, sec. 3.2; @thanh01].

% Code blocks are used to include source code listings into the
% document:

%     #include <stdio.h>
%     #include <unistd.h>
%     #include <sys/types.h>
%     #include <sys/wait.h>
%     // This is a comment
%     int main(int argc, char **argv)
%     {
%         while (--c > 1 && !fork());
%         sleep(c = atoi(v[c]));
%         printf("%d\n", c);
%         wait(0);
%         return 0;
%     }

% There is an alternative syntax for code blocks that allows you to
% specify additional information, such as the language of the source
% code. This information can be used for syntax highlighting:

% ``` sh
% #!/bin/sh
% fac() {
%   if [ "$1" -leq 1 ]; then
%     echo 1
%   else
%     echo $(("$1" * fac $(("$1" - 1))))
%   fi
% }
% ``````````````

% ~~~~~~ Ruby
% # Here's a way to empty an array.
% joe = [ 'eggs.', 'some', 'break', 'to', 'Have' ]
% print(joe.pop, " ") while joe.size > 0
% print "\n"
% ~~~~~~

% \end{markdown*}
% \shorthandon{-}

% \chapter{Inserting the bibliography}
% After linking a bibliography data\-base files to the document using
% the \verb"\"\texttt{thesis\discretionary{-}{}{}setup\{bib\discretionary{=}{=}{=}

% \{\textit{file1},\textit{file2},\,\ldots\,\}\}} command, you can
% start citing the entries. This is just dummy text
% \parencite{borgman03} lightly sprinkled with citations
% \parencite[p.~123]{greenberg98}. Several sources can be cited at
% once: \cite{borgman03,greenberg98,thanh01}.
% \citetitle{greenberg98} was written by \citeauthor{greenberg98} in
% \citeyear{greenberg98}. We can also produce \textcite{greenberg98}%
% \ or %% Let us define a compound command:
% \def\citeauthoryear#1{(\textcite{#1},~\citeyear{#1})}%
% \citeauthoryear{greenberg98}%
% . The full bibliographic citation is:
% \emph{\fullcite{greenberg98}}. We can easily insert a bibliographic
% citation into the footnote\footfullcite{greenberg98}.

% The \verb"\nocite" command will not generate any
% output\nocite{muni}, but it will insert its arguments into
% the bibliography. The \verb"\nocite{*}" command will insert all the
% records in the bibliography database file into the bibliography.
% Try uncommenting the command
% %% \nocite{*}
% and watch the bibliography section come apart at the seams.

% When typesetting the document for the first time, citing a
% \texttt{work} will expand to [\textbf{work}] and the
% \verb"\printbibliography" command will produce no output. It is now
% necessary to generate the bibliography by running \texttt{biber
% \jobname.bcf} from the command line and then by typesetting the
% document again twice. During the first run, the bibliography
% section and the citations will be typeset, and in the second run,
% the bibliography section will appear in the table of contents.

% The \texttt{biber} command needs to be executed from within the
% directory, where the \LaTeX\ source file is located. In Windows,
% the command line can be opened in a directory by holding down the
% \textsf{Shift} key and by clicking the right mouse button while
% hovering the cursor over a directory.  Select the \textsf{Open
% Command Window Here} option in the context menu that opens shortly
% afterwards.

% With online services -- such as Overleaf -- or when using an
% automatic tool -- such as \LaTeX MK -- all commands are executed
% automatically. When you omit the \verb"\printbibliography" command,
% its location will be decided by the template.


% \chapter{Inserting the index}
% After using the \verb"\makeindex" macro and loading the
% \texttt{makeidx} package that provides additional indexing
% commands, index entries can be created by issuing the \verb"\index"
% command. \index{dummy text|(}It is possible to create ranged index
% entries, which will encompass a span of text.\index{dummy text|)}
% To insert complex typographic material -- such as $\alpha$
% \index{alpha@$\alpha$} or \TeX{} \index{TeX@\TeX} --
% into the index, you need to specify a text string, which will
% determine how the entry will be sorted. It is also possible to
% create hierarchal entries. \index{vehicles!trucks}
% \index{vehicles!speed cars}

% After typesetting the document, it is necessary to generate the
% index by running
% \begin{center}%
%   \texttt{texindy -I latex -C utf8 -L }$\langle$\textit{locale}%
%   $\rangle$\texttt{ \jobname.idx}
% \end{center}
% from the command line, where $\langle$\textit{locale}$\rangle$
% corresponds to the main locale of your thesis -- such as
% \texttt{english}, and then typesetting the document again.

% The \texttt{texindy} command needs to be executed from within the
% directory, where the \LaTeX\ source file is located. In Windows,
% the command line can be opened in a directory by holding down the
% \textsf{Shift} key and by clicking the right mouse button while
% hovering the cursor over a directory. Select the \textsf{Open Command
% Window Here} option in the context menu that opens shortly
% afterwards.

% With online services -- such as Overleaf -- the commands are
% executed automatically, although the locale may be erroneously
% detected, or the \texttt{makeindex} tool (which is only able to
% sort entries that contain digits and letters of the English
% alphabet) may be used instead of \texttt{texindy}. In either case,
% the index will be ill-sorted.


